\begin{table}[H]
  \centering
  \small
  \caption{Description textuelle du cas d’utilisation : «~Créer un site~»}
  \label{tab:cu-gerer-sites}
  \PFETableSetup
  \PFETableRowColors
  \PFETableArrayStretch
  \PFETableTabularXBegin

  \begin{tabularx}{\PFETableBlockWidth}{@{}>{\raggedright\arraybackslash}p{3.2cm} >{\raggedright\arraybackslash}X@{}}
    
    \tableHeaderStyle 
    \tableHeaderCell{Champ} & \tableHeaderCell{Contenu} \\

    \textbf{Titre} & Créer un site \\

    \textbf{Acteurs} & Utilisateur Corporate \\

    \textbf{Objectif} &
    Permettre à l’utilisateur Corporate de créer les sites de Coficab afin d’organiser les activités CSR et structurer l’accès aux données par site. \\

    \textbf{Préconditions} &
    1. L’utilisateur Corporate est authentifié.\par
    2. Un token est obtenu pour accéder aux ressources protégées. \\

    \textbf{Postconditions} &
    1. Les informations du site sont enregistrées.\par
    2. Le nouveau site est disponible pour les autres modules de la plateforme. \\

    \textbf{Scénario nominal} &
    1. L’utilisateur Corporate accède au module de gestion des sites. \par
    2. Le système affiche la liste des sites existants. \par
    3. L’utilisateur choisit l’action : créer un site. \par
    4. L’utilisateur saisit les informations du site (nom, code, pays, région…). \par
    5. Le système valide les données saisies. \par
    6. Si les données sont valides, le système enregistre la création du site et met à jour la liste des sites. \\

    \textbf{Scénario alternatif} &
    \textbf{A. Site déjà existant :} \par
    A.1 L’utilisateur tente de créer un site avec un code déjà existant. \par
    A.2 Le système rejette l’opération. \par
    A.3 Un message est affiché : «~Code site déjà existant~». \\

  \end{tabularx}%

  \PFETableTabularXEnd
\end{table}
